\section{Introduction}\label{sec:intro}
Deep learning models improve with the amount of labelled training data; in computer vision, performance has been shown to grow with the logarithm of the training-set size.\cite{sun2017} Medical imaging rarely offers data at this scale: the main obstacle is the availability of sufficiently large, curated and representative datasets with expert labels.\cite{willemink2020} The shortage is most severe when the label is a clinical outcome rather than an image finding, because outcomes can only be collected by following patients over time.

This is the case for overall survival prediction from pretreatment computed tomography (CT) in non-small-cell lung cancer (NSCLC), where most prognostic studies are retrospective and single-centre, with cohorts that usually include only a few hundred patients, of whom relatively few have complete survival follow-up.\cite{lun2026,qin2026} Public datasets do not change this picture: LUNG1, one of the most widely used CT survival cohorts, includes 422 patients,\cite{aerts2014} and deep models for lung cancer outcomes have been developed on as few as 130.\cite{verma2024}

For cohorts of this size, foundation models offer an alternative to training a deep network from scratch. They are pretrained on broad data at scale and are intended to be adapted to many downstream tasks with very little task-specific labelled data.\cite{bommasani2021, moor2023} For a small survival cohort, this adaptation can range from fitting a survival model on the features of a frozen encoder to adding parameter-efficient modules such as low-rank adaptation (LoRA)\cite{hu2021} or fine-tuning part of the encoder. Several foundation models now exist for CT, trained with different objectives: language--image alignment (CT-CLIP and COLIPRI),\cite{hamamci2026ctclip,wald2025colipri} masked autoencoding (TANGERINE),\cite{mcconnell2026tangerine} joint-embedding prediction (Rad-JEPA-3D),\cite{trinh2026radjepa} and contrastive pretraining on lesion patches.\cite{pai2024} General foundations trained on natural images, such as DINOv3, can also be applied to CT slices.\cite{simeoni2025dinov3}

For overall survival in NSCLC, these foundation models join a range of earlier approaches: clinical models of stage and tumour volume, handcrafted radiomic features of the segmented tumour,\cite{aerts2014} and supervised deep networks trained on tumour patches.\cite{hosny2018} Frozen features from lesion-centred foundation models have already been evaluated for this task,\cite{pai2024,aerts2025tib} whereas whole-volume CT foundation models have been applied to it only recently.\cite{neukirch2026,allgower2026} Each of these approaches has been evaluated in its own setting, with different cohorts, endpoints and baselines, so it remains unclear how they compare with one another.

The few studies that compared several of them used binary two-year endpoints, lesion-centred inputs or a single foundation model (Section~\ref{sec:related}).\cite{aerts2025tib,neukirch2026,allgower2026} A comparison within one protocol, with the same cohorts, folds, continuous-time endpoint and clinical baselines, would also allow testing the premise that motivates foundation models in this setting: if they compensate for small cohorts, adapting them to survival should improve discrimination, and their advantage over simpler models should grow as training patients become fewer. It would further show whether whole-volume embeddings draw on the tumour or on the surrounding anatomy, and whether their predictions transfer to other populations in absolute risk as well as in ranking, since external evaluations have so far reported discrimination only.

We therefore compared five foundation models (COLIPRI, TANGERINE, Rad-JEPA-3D, CT-CLIP and DINOv3), two deep networks trained on the tumour region, a radiomic model and clinical models within a single protocol on LUNG1, with external evaluation on NSCLC-Radiogenomics,\cite{bakr2018} and addressed four research questions:
\begin{enumerate}
\item[RQ1] How much prognostic information is accessible from frozen foundation-model representations, relative to clinical, radiomic and tumour-centred models?
\item[RQ2] How do trainable survival heads, adaptation of the encoder and the training-label budget change survival discrimination?
\item[RQ3] How sensitive are the representations and their survival predictions to the tumour?
\item[RQ4] Do discrimination, incremental clinical value and calibration transfer to an external cohort?
\end{enumerate}

In answering these questions within a single protocol, this study makes four contributions, which correspond in order to RQ1--RQ4:
\begin{itemize}
\item A unified survival benchmark of CT and general-purpose foundation models, tumour-centred comparators, radiomics and clinical baselines on the same cohorts, folds, endpoint and external evaluation, including encoders not previously evaluated for survival (COLIPRI, TANGERINE, Rad-JEPA-3D and DINOv3).
\item A comparison of linear Cox models on frozen embeddings, trainable survival heads, LoRA and partial fine-tuning, together with learning curves across training-label budgets, on identical development and test patients.
\item An analysis of tumour dependence that masks tumour tokens inside each frozen encoder under matched spatial controls, separating changes in the representation from changes in discrimination.
\item An external evaluation that reports calibration alongside discrimination, separating the transfer of risk ranking from the transfer of absolute risk.
\end{itemize}

The remainder of the paper is organized as follows. Section~\ref{sec:related} reviews tumour-centred prognostic models, CT foundation models and prior benchmarks of their survival performance. Section~\ref{sec:methods} describes the cohorts, representations, survival models, adaptation methods, perturbations and external evaluation. Section~\ref{sec:results} reports the results for RQ1--RQ4, Section~\ref{sec:discussion} discusses their interpretation, Section~\ref{sec:limitations} the limitations of the study, and Section~\ref{sec:conclusions} concludes.

\gap{Add the GitHub repository URL once the code is public.}
